\section{总结与延伸}

\subsection{总结表}
以下表格汇总本讲的主要模块、核心概念与实践重点，为复盘与 audit 提供直接的检索路径：

\begin{table}[H]
\centering
\small
\begin{tabular}{p{0.27\textwidth}p{0.33\textwidth}p{0.33\textwidth}}
\toprule
模块 & 核心概念 & 工程/实践要点 \\
\midrule
模仿\(\to\)强化 & Online RL + 采样闭环 & 策略参数直接控制数据分布 \\
策略梯度推导 & log-gradient trick & 只优化 \(\pi_\theta\)，无需懂动力学 \\
直觉与改进 & reward-to-go + baseline & advantage 估计降低 variance \\
稳定性技术 & GAE + Entropy + Trust Region & PPO/TRPO + entropy bonus \\
实现细节 & 超参 + 调试指标 & monitoring reward curve 和 entropy \\
Off-policy & Importance weights & 仅在策略近似时复用旧数据 \\
\bottomrule
\end{tabular}
\caption{按模块划分的关键输出与风险控制}
\end{table}

\subsection{延伸阅读}
\begin{itemize}
    \item Williams, \textit{Simple Statistical Gradient-Following Algorithms for Connectionist Reinforcement Learning} (REINFORCE 原文, 1992)。
    \item Schulman et al., \textit{High-Dimensional Continuous Control Using Generalized Advantage Estimation}。
    \item Sergey Levine, \textit{CS285 Lecture on Policy Gradients}：另一套策略梯度推导与工程技巧。
    \item Schulman et al., \textit{Proximal Policy Optimization Algorithms}：理解 trust region/clipped objective 的落地方法。
    \item Sutton \& Barto, \textit{Reinforcement Learning: An Introduction}（第13章）涵盖 policy gradient 与 actor-critic。
\end{itemize}

\subsection{本章小结}
本讲围绕策略梯度的形式推导、直觉解释与工程实践展开，强调 variance control、entropy regularization 以及实现层面的 logging/checkpoint 策略。Off-policy 的讨论指出了该系列方法的数据瓶颈，为下一步引入 actor-critic 与 value-based hybrid 亮出舞台。

\end{document}
